\subsection{Eberlein 定理与 Šmulian 定理}

定理 1 (Eberlein). --- 设 E 是一个 Hausdorff 且拟完备的局部凸空间，T 是 E 上与 E 和 E′ 之间的对偶相容的一个拓扑，A 是 E 的一个子集。为使 A 对于 T 相对紧，必须且只需 A 中点的每个无穷序列对于 T 在 E 中有一个极限点。

所述条件显然是必要的。

现在假设 A 中点的每个无穷序列对于 T 有一个极限点，从而对于较粗的拓扑 \(\sigma(\mathrm{E}, \mathrm{E}')\) 也有一个极限点。那么 A 对于 T 预紧 (IV, 第 32 页, 命题 1)；为使 A 对于 T 相对紧，必须且只需它对于 \(\sigma(\mathrm{E}, \mathrm{E}')\) 相对紧（出处同上）。因此只需对 T 是弱化拓扑 \(\sigma(\mathrm{E}, \mathrm{E}')\) 的情形证明本定理。

设 \(\hat{\mathrm{E}}\) 是 E 的完备化，我们照例把它与代数对偶 \(\mathrm{E}^{\prime *}\)（\(\mathrm{E}'\) 的代数对偶）的一个子空间视为同一 (III, 第 21 页, 定理 2)。设 \(\mathrm{E}_{\sigma}\)、\(\hat{\mathrm{E}}_{\sigma}\) 与 \(\mathrm{E}_{\sigma}^{\prime *}\) 分别表示空间 E、\(\hat{\mathrm{E}}\) 与 \(\mathrm{E}^{\prime *}\)，各赋以拓扑 \(\sigma (\mathrm{E},\mathrm{E}')\)、\(\sigma (\hat{\mathrm{E}},\mathrm{E}')\) 与 \(\sigma (\mathrm{E}^{\prime *},\mathrm{E}^{\prime})\)。

设 \((x_{i}^{\prime})_{i\in\mathbb{I}}\) 是向量空间 \(E^{\prime}\) 在域 K 上的一个基。映射 \(f\mapsto(f(x_{i}^{\prime}))_{i\in\mathbb{I}}\) 是一个同胚 \(\phi\)（从 \(E_{\sigma}^{\prime*}\) 到 \(K^{I}\) 上）；对每个 \(i\in I\)，A 在映射 \(x_{i}^{\prime}\)（从 E 到 K 中的映射）下的像是相对紧的：因为 K 可度量化，且 \(x_{i}^{\prime}(A)\) 的元素的每个无穷序列有一个极限点。由此推出 \(\phi(A)\) 在 \(K^{I}\) 中相对紧，从而闭包 \(\overline{A}\)（A 在 \(E_{\sigma}^{\prime*}\) 中的闭包）是紧的。

其次我们证明 \(\overline{A}\) 包含在 \(\hat{E}\) 中。设 H 是 \(E'\) 的一个等度连续子集；设 X 是它对于 \(\sigma(E', E)\) 的闭包；X 是紧的 (III, 第 17 页, 推论 2)。对每个 \(x \in E'^{*}\)，设 \(\phi_{x}\) 是 \(x' \mapsto \langle x, x' \rangle\) 在 X 上的限制；设 \(\tilde{A} \subset \mathcal{C}_{s}(X)\) 是当 x 取遍 A 时函数 \(\phi_{x}\) 所成的集。鉴于关于 A 的假设，\(\tilde{A}\) 的元素的每个无穷序列在 \(\mathcal{C}_{s}(X)\) 中有一个极限点；由命题 2 (IV, 第 33 页)，集 \(\tilde{A}\) 因而在 \(\mathcal{C}_{s}(X)\) 中相对紧。由此推出，对每个 \(a \in \overline{A}\)，X 上的函数 \(\phi_{a}\) 是连续的。包含关系 \(\overline{A} \subset \hat{E}\) 于是由 III, 第 21 页 的定理 2 得出。

现在我们证明 \(\overline{\mathbf{A}}\) 包含在 E 中。由于 A 在 \(\mathrm{E}_{\sigma}\) 中预紧 (IV, 第 32 页, 命题 1)，它在 \(\mathrm{E}_{\sigma}\) 中有界 (III, 第 3 页, 命题 2)，从而在 E 中也有界 (IV, 第 1 页, 命题 1)。设 C 是 A 在 E 中的闭均衡凸包。那么 C 有界（因 A 有界），从而完备（因 E 拟完备）。换句话说，C 是 \(\hat{\mathrm{E}}\) 的一个凸闭子集，因而也是 \(\hat{\mathrm{E}}_{\sigma}\) 的凸闭子集 (IV, 第 1 页, 命题 1)。由于 \(\mathbf{A} \subset \mathbf{C}\) 且 \(\hat{\mathrm{E}}_{\sigma}\) 的拓扑由 \(\mathrm{E}_{\sigma}^{\prime *}\) 的拓扑诱导，我们有 \(\overline{\mathbf{A}} \subset \mathbf{C}\)，从而 \(\overline{\mathbf{A}} \subset \mathbf{E}\)。

由于 \(E_{\sigma}\) 的拓扑由 \(E_{\sigma}^{\prime*}\) 的拓扑诱导，子集 \(\overline{A}\)（\(E_{\sigma}\) 的子集）是紧的，定理 1 随之得证。

定理 2 (Šmulian). --- 设 E 是一个 Fréchet 空间，A 是 E 的一个子集。设 \(E_{\sigma}\) 表示赋以弱化拓扑的空间 E。下列条件等价：

\begin{enumerate}
\def\labelenumi{(\roman{enumi})}
\item
  A 在 \(\mathrm{E}_{\sigma}\) 中相对紧；
\item
  A 中点的每个无穷序列在 \(\mathrm{E}_{\sigma}\) 中有一个极限点；
\item
  从 \(\mathbf{A}\) 中点的每个无穷序列，可以抽取一个在 \(\mathrm{E}_{\sigma}\) 中收敛的序列。
\end{enumerate}

(i) 与 (ii) 的等价性由 Eberlein 定理得出，且 (iii) 显然蕴含 (ii)。

我们来证明 (i) 蕴含 (iii)。假设 A 在 \(E_{\sigma}\) 中的闭包 B 是紧的，且 \((x_{n})_{n\in\mathbf{N}}\) 是 A 的点的一个序列。设 F 是 E 中包含诸 \(x_{n}\) 的最小闭向量子空间，它是一个满足第一可数公理的 Fréchet 空间。由于 F 在 \(E_{\sigma}\) 中闭，且 F 上的拓扑 \(\sigma(\mathrm{F},\mathrm{F}')\) 由 \(\sigma(\mathrm{E},\mathrm{E}')\) 诱导，集 \(B\cap F\) 对于 \(\sigma(\mathrm{F},\mathrm{F}')\) 是紧的。根据 IV, 第 32 页 的注记，从 \((x_{n})_{n\in\mathbf{N}}\) 中抽取一个对于 \(\sigma(\mathrm{E},\mathrm{E}')\)（或者等价地，对于 \(\sigma(\mathrm{F},\mathrm{F}')\)）收敛的序列，其存在性是下述引理的推论：

引理 1. --- 设 F 是一个满足第一可数公理的 Fréchet 空间。F 的每个对于拓扑 T（由 \(\sigma(\mathrm{F},\mathrm{F}')\) 诱导）紧的子集 C 对于该拓扑可度量化。

由于 \(F'\) 上的预紧收敛拓扑比拓扑 \(\sigma(F', F)\) 细，在 \(F_{s}'\) 中存在一个处处稠密的可数子集 (III, 第 18 页, 推论 1)。于是集 C 可以与 \(K^{D}\) 的一个子集视为同一，且 \(K^{D}\) 的拓扑（它可度量化 (GT, IX, § 2, No.~8)）在 C 上诱导的拓扑比由 \(\sigma(F, F')\) 诱导的拓扑粗，而 C 对于后者是紧的。因此这两个拓扑相同 (GT, I, § 9, No.~4, 推论 3)。

Šmulian 定理可以推广到 E 是一列 Fréchet 空间的严格归纳极限的情形 (IV, 第 67 页, 习题 2)。

